% !TEX root = ../aa_SoM_Chapters.tex

% ö = \%c3\%b6
% ä = \%C3\%A4

\xdef\sectiontitle{\Isectiontitle} %aiheuttama

%%%%%%%%%%%%%%%
\begin{frame}[label=1_9_a]%[label=1_9_TOC1]
\frametitle{\texorpdfstring{\culine[\sectiontitleulinecolor]{\sectiontitle}}{\sectiontitle} }

\vspace{0.5cm}

\begin{itemize}[leftmargin=0.5cm,itemsep=2mm]
    \item[1.1]<1-> \hyperlink{1_1_a}{\IaSubsectiontitle}
    \item[1.2]<1-> \hyperlink{1_2_a}{\IbSubsectiontitle}	
    \item[1.3]<1-> \hyperlink{1_3_a}{\IcSubsectiontitle}	
    \item[1.4]<1-> \hyperlink{1_4_a}{\IdSubsectiontitle}
    \item[1.5]<1-> \hyperlink{1_5_a}{\IeSubsectiontitle}
    \item[1.6]<1-> \hyperlink{1_6_a}{\IfSubsectiontitle}
    \item[1.7]<1-> \hyperlink{1_7_a}{\IgSubsectiontitle}
\end{itemize}

\vspace{-1cm}
\begin{itemize}[leftmargin=9.5cm,itemsep=2mm]
    \item[1.8]<1-> \hyperlink{1_8_a}{\IhSubsectiontitle}
	\item[1.9]<1-> \hyperlink{1_9_a}{\CDAlert[blue]{\IiSubsectiontitle}}
	\item[1.10]<1-> \hyperlink{1_10_a}{\IjSubsectiontitle}
	\item[1.11]<1-> \hyperlink{1_11_a}{\IkSubsectiontitle}
\end{itemize}

%\endofslide 
\end{frame}
%%%%%%%%%%%%%%%


\xdef\sectiontitle{\IiSubsectiontitle}

%%%%%%%%%%%%%%%
\begin{frame}%[label=1_9_a]
\frametitle{\texorpdfstring{\culine[\sectiontitleulinecolor]{\sectiontitle}}{\sectiontitle} }

%\vspace{1.5mm}
\begin{minipage}{8.1cm}
\begin{itemize}[itemsep=1.7mm]
	\item Lasers are a good example of a photon gas system
	\item They are devices where high energy density of electromagnetic radiation \appear{is introduced within a cavity}\appear{, outputting a \Alert[red]{coherent unidirectional monochromatic beam} of light}

\item The basic components are:
\begin{itemize}
\small
	\item A \CDAlert[fuchsia]{three-level system} \appear{in which atoms in the upper level "2"} \appear{make a transition to the lower level "1"} \appear{that is stimulated by the presence of photons of the right frequency}
\item An \CDAlert[blue]{optical cavity} \appear{to provide an environment with the right kind of photons}
\item \CDAlert[red]{Pumping mechanism} to re-populate the upper level "2" \appear{for repeated operation}
\end{itemize}
\end{itemize}
\end{minipage}

\xdef\loc{east}
\begin{tikzpicture}[remember picture,overlay] % anchor=south
    \node (A) [yshift=0cm,xshift=\figXshift, anchor=\loc] at (current page.\loc) {\includegraphics[page=1,height=8cm]{Figures_booklet/SOM_9_Statistical_mechanics_figures/SOM_Figure_21.png}}; %scale=\figscale. % 8cm max hei
\end{tikzpicture}

\endofslide 
%\endofsection
\end{frame}
%%%%%%%%%%%%%%%

%%%%%%%%%%%%%%%
\begin{frame}
\frametitle{\texorpdfstring{\culine[\sectiontitleulinecolor]{\sectiontitle}}{\sectiontitle} }

\begin{itemize}
	\item Practical operating principles of laser:

\item[] \begin{itemize}[itemsep=2.3mm] \small
	\item Parallel mirrors are placed at the medium's two ends
  \item Photons not parallel the axis \appear{have relatively few opportunities to induce stimulated emission}
  \item \Alert[blue]{Photons parallel to axis reflect back and forth}\appear{, providing many opportunities}
  \item The length $L$ \appear{is tuned to standing wave:} \appear{$L=n \lambda / 2$}
    % 6.5 cm = midway, 10 cm = 3/4 
%\vspace{2.5mm}
\end{itemize}
\end{itemize}
\begin{minipage}{6.25cm}
\begin{itemize}
\item[] \begin{itemize}[itemsep=2.3mm]  \small
  \item A resonance is formed inside this optical cavity
  \item \CDAlert[blue]{Parallel radiation gets amplified}
  \item Off-axis radiation diminishes
  \item \Alert[red]{Laser wavelength becomes narrower than natural line width}
  \item More energy can be fed into the laser by an 'optical \Alert[fuchsia]{pumping process}'
\end{itemize}
\end{itemize}
  \end{minipage}

\xdef\loc{south east}
\begin{tikzpicture}[remember picture,overlay] % anchor=south
    \node (A) [yshift=-0.25*\figYshift,xshift=\figXshift, anchor=\loc] at (current page.\loc) {\includegraphics[page=1,width=6.7cm]{Figures_booklet/SOM_9_Statistical_mechanics_figures/SOM_Figure_23.png}}; %scale=\figscale. % 8cm max hei
\end{tikzpicture}


\endofslide 
%\endofsection
\end{frame}
%%%%%%%%%%%%%%%

%%%%%%%%%%%%%%%
\begin{frame}
\frametitle{\texorpdfstring{\culine[\sectiontitleulinecolor]{\sectiontitle}}{\sectiontitle} }

\begin{minipage}{7cm}
\begin{itemize}[itemsep=2.7mm]
	% 6.5 cm = midway, 10 cm = 3/4 
\vspace{2mm}
\small
	\item \CDAlert[red]{Spontaneous emission} \appear{occurs when an electron in a higher atomic energy level} \appear{drops to lower level in a short but unpredictable time}\appear{, emitting a photon}

\item \CDAlert[blue]{Absorption} \appear{occurs when an electron is promoted to higher atomic energy level} \appear{by absorbing a photon}

\item In \CDAlert[fuchsia]{stimulated emission}\appear{, an electron in higher energy level is influenced by a passing photon}\appear{, whose energy is precisely the difference between the upper and lower level} 
 \item[] \begin{itemize}
	\item[=>] As a result, another photon with identical energy \appear{and wavelength} \appear{to the first photon is emitted}\appear{, and moreover in the same direction and with a \CDAlert[fuchsia]{correlated phase} with the first one}
\end{itemize}
 

\end{itemize}
\end{minipage}

\xdef\loc{east}
\begin{tikzpicture}[remember picture,overlay] % anchor=south
    \node (A) [yshift=0cm,xshift=\figXshift, anchor=\loc] at (current page.\loc) {\includegraphics[page=1,height=8cm,trim={0 0 13cm 0},clip]{Figures_booklet/SOM_9_Statistical_mechanics_figures/SOM_Figure_22.png}}; %scale=\figscale. % 8cm max hei
\end{tikzpicture}


\endofslide 
%\endofsection
\end{frame}
%%%%%%%%%%%%%%%

%%%%%%%%%%%%%%%
\begin{frame}
\frametitle{\texorpdfstring{\culine[\sectiontitleulinecolor]{\sectiontitle}}{\sectiontitle} }

\begin{minipage}{7cm}
\begin{itemize}\small 
	% 6.5 cm = midway, 10 cm = 3/4 
\vspace{2mm}
	\item How to describe these three processes mathematically?
For simplicity, gases are studied (simpler energy state treatment than in solids)

\item In spontaneous emission\appear{, the \CDAlert[red]{rate $R_{\text{spon}}$}} \appear{must be proportional to the \Alert[red]{number of atoms in the excited state, $N_{2}$}}\appear{, the proportionality constant being $A_{\text{spon}}$} \appear{(\Alert[red]{$R_{\text{spon}}=A_{\text {spon}} N_2$})}

\item \CDAlert[blue]{Rate of absorption} must be proportional to the number of particles in the lower level, $N_{1}$\appear{, but also proportional to the number of photons of the right energy which are entering the atom $Y(\Delta E)$}\appear{, where $\Delta E=E_{2}-E_{1}$.} \appear{The \Alert[blue]{proportionality constant is $B_{abs}$}} \appear{(\Alert[blue]{$R_{abs}=B_{abs} N_{1} Y(\Delta E)$})}

\item Also for \CDAlert[fuchsia]{stimulated emission}\appear{, we need the number of available atoms in the higher level}\appear{, as well as the number of photons of right energy}\appear{, which can trigger the emission} \appear{(\Alert[fuchsia]{$R_{\text{stim}}=B_{\text{stim}} N_{2} Y(\Delta E)$})}
\end{itemize}
\end{minipage}

\xdef\loc{east}
\begin{tikzpicture}[remember picture,overlay] % anchor=south
    \node (A) [yshift=0cm,xshift=\figXshift, anchor=\loc] at (current page.\loc) {\includegraphics[page=1,height=8cm,trim={0 0 13cm 0},clip]{Figures_booklet/SOM_9_Statistical_mechanics_figures/SOM_Figure_22.png}}; %scale=\figscale. % 8cm max hei
\end{tikzpicture}


\endofslide 
%\endofsection
\end{frame}
%%%%%%%%%%%%%%%

%%%%%%%%%%%%%%%
\begin{frame}
\frametitle{\texorpdfstring{\culine[\sectiontitleulinecolor]{\sectiontitle}}{\sectiontitle} }

\begin{itemize}
	\item The coefficients $A_{\text {spon}}$\appear{, $B_{abs}$} \appear{and $B_{\text {stim}}$} \appear{are known as the \CDAlert[red]{Einstein's coefficients} }
	
	\item In an equilibrium, these processes are in a balance:
\[
\begin{aligned}
& & R_{\text {spon}} \appear{+R_{\text {stim}}}  &\equal \appear{R_{abs}} \\
 &\appear{\Rightarrow}& \qquad \appear{A_{\text {spon}} N_{2}} \appear{+B_{\text {stim}} N_{2} Y(\Delta E)} & \equal \appear{B_{abs} N_{1} Y(\Delta E)} \qquad \quad 
\end{aligned}
\]

\item From where we solve for $Y(\Delta E)$\appear{, and using $\appear{N_{1} / N_{2}}  \equal \appear{e^{-\frac{E_{1}}{k_{B} T}}} \appear{/ e^{-\frac{E_{2}}{k_{B} T}}}$}\appear{, we get:}
\[
\begin{aligned}
&\Rightarrow& 
\qquad \appear{Y(\Delta E)} & \equal \frac{A_{\text {spon}} \appear{/ B_{abs}}}{\frac{N_{1}}{N_{2}}  \minus \frac{B_{\text {stim}}}{B_{abs}}}
 \equal 
 \frac{A_{\text {spon}} / B_{abs}}{\frac{e^{-E_{1} / k_{B} T}}{e^{-E_{2} / k_{B} T}}  \minus \frac{B_{\text{stim}}}{B_{abs}}} \\
 \vspace{6mm}
&\Rightarrow&
 \appear{Y(\Delta E)} & \equal \frac{A_{\text {spon}} / B_{abs}}{\appear{e^{\Delta E / k_{B} T}}  \minus \frac{B_{\text {stim}}}{B_{abs}}}
  \equal 
  \frac{A_{\text {spon}} / B_{\text {stim}}}{\Frac{B_{abs}}{B_{\text{stim}}} e^{\Delta E / k_{B} T}-1}  \qquad \quad 
\end{aligned}
\]
\end{itemize}

\endofslide 
%\endofsection
\end{frame}
%%%%%%%%%%%%%%%

%%%%%%%%%%%%%%%
\begin{frame}
\frametitle{\texorpdfstring{\culine[\sectiontitleulinecolor]{\sectiontitle}}{\sectiontitle} }

\begin{itemize}
	\item The quantity $Y(\Delta E)$ is not strictly speaking a number of photons entering the atom\appear{, but rather a photon gas spectral energy density for the quanta of right energy}\appear{, given by \CDAlert[red]{Planck's law}} 
	
	\item So, we can compare:
\[
\begin{aligned}
Y(\Delta E) &\equal \frac{A_{\text {spon}} / B_{abs}}{\frac{N_{1}}{N_{2}} \minus \frac{B_{\text {stim}}}{B_{abs}}}
 \equal 
 \frac{A_{\text{spon}} / B_{abs}}{\frac{e^{-E_{1} / k_{B} T}}{e^{-E_{2} / k_{B} T}}-\frac{B_{\text {stim}}}{B_{abs}}}
  \equal 
  \frac{A_{\text {spon}} / B_{abs}}{e^{\Delta E / k_{B} T}  \minus \frac{B_{\text {stim}}}{B_{\text {abs }}}}
   \equal 
   \frac{A_{\text {spon}} / B_{\text {stim}}}{\Frac{B_{abs}}{B_{\text {stim}}} e^{\Delta E / k_{B} T}-1} \\
\appear{d U_{photon}(\Delta E)} &\equal 
 \frac{h f^{3}}{e^{\Frac{\Delta E}{k_{B} T}}-1} \frac{8 \pi V}{c^{3}} \appear{d f} \quad \Rightarrow \quad \frac{1}{V} \frac{d U_{photon}}{d f}
  \equal \frac{8 \pi}{c^{3}} \frac{h f^{3}}{e^{\Frac{\Delta E}{k_{B} T}}-1}  \equal  \appear{Y(\Delta E)}
\end{aligned}
\]
\item Therefore, \CDAlert[blue]{in equilibrium} the below holds:
\[
B_{\text {stim}}=B_{abs} \qquad \appear{\text{resulting~in}} \qquad \appear{\frac{A_{\text {spon}}}{B_{\text {stim}}}  \equal \frac{8 \pi h f^{3}}{c^{3}}}
\]
\end{itemize}

%\xdef\loc{east}
%\begin{tikzpicture}[remember picture,overlay] % anchor=south
%    \node (A) [yshift=0cm,xshift=\figXshift, anchor=\loc] at (current page.\loc) {\includegraphics[page=1,height=7cm]{Figures_booklet/SOM_9_Statistical_mechanics_figures/SOM_Figure_17.png}}; %scale=\figscale. % 8cm max hei
%\end{tikzpicture}


\endofslide 
%\endofsection
\end{frame}
%%%%%%%%%%%%%%%

%%%%%%%%%%%%%%%
\begin{frame}
\frametitle{\texorpdfstring{\culine[\sectiontitleulinecolor]{\sectiontitle}}{\sectiontitle} }

\begin{itemize}
	\item Note, that condition \appear{$B_{\text{stim}}=B_{abs}$} \appear{means that an incoming photon is equally probable to trigger an absorption} \appear{as a stimulated emission}

\item Frankly speaking, for lasers we do not favour this\appear{, we wish for the stimulation to dominate}\appear{, since we wish to have a strong coherent beam}\appear{, made by \CDAlert[red]{utilizing stimulated emission}}
\end{itemize}

\xdef\loc{south}
\begin{tikzpicture}[remember picture,overlay] % anchor=south
    \node (A) [yshift=-0.25*\figYshift,xshift=0*\figXshift, anchor=\loc] at (current page.\loc) {\includegraphics[page=11,height=5.3cm]{Figures_booklet/SOM_9_Statistical_mechanics_figures/Tikz_SOM_Statistical_figures.pdf}}; %scale=\figscale. % 8cm max height
\end{tikzpicture}

\endofslide 
%\endofsection
\end{frame}
%%%%%%%%%%%%%%%


%%%%%%%%%%%%%%%
\begin{frame}
\frametitle{\texorpdfstring{\culine[\sectiontitleulinecolor]{\sectiontitle}}{\sectiontitle} }

\begin{itemize}
	\item At normal temperature\appear{, in equilibrium}\appear{, nearly all of the atoms are in the ground state} 
	\item Thus $N_{1} \gg N_{2}$\appear{, and for actual rates}\appear{, it follows that in normal temperature $N_{1} B_{abs} \gg N_{2} B_{\text{stim}}$}  
	\implies{ \Alert[red]{stimulated emission is not actually pronounced} }

\item Using optical pumping\appear{, i.\,e. an additional light source  with right frequency}\appear{, the ground states can be excited into higher states} 
\implies{ Population of a higher level with energy $E_{3}$ increases }

\item If the medium of the laser \appear{(often a \CDAlert[tuniblue]{transparent solid})} \appear{has a \CDAlert[fuchsia]{metastable state $E_{2}$}}\appear{, then many of the upper level $E_{3}$ states can drop down to this metastable state $E_{2}$} 
\implies{ Possible to create a situation where $N_{1}<N_{2}$ }

\end{itemize}

\endofslide 
%\endofsection
\end{frame}
%%%%%%%%%%%%%%%

%%%%%%%%%%%%%%%
\begin{frame}
\frametitle{\texorpdfstring{\culine[\sectiontitleulinecolor]{\sectiontitle}}{\sectiontitle} }

\begin{itemize}
	\item First, one uses optical pumping \appear{so that $N_{1} \approx \appear{N_{2}} \approx \appear{N_{3}}$} 
	\item When the pumping is stopped\appear{, very soon afterwards}\appear{, the level $E_{3}$ will depopulate}\appear{, and soon there will be a high population of $E_{2}$}\appear{, and also $N_{1}<N_{2}$}
	
	\item This phenomenon\appear{/condition} \appear{is called \CDAlert[red]{population inversion}} 

\end{itemize}
\vspace{2.5mm}
\begin{minipage}{3.8cm}
\begin{itemize}
	 \item As a result, once the photon gas in the laser starts
	 to trigger more of the $E_{2}$
	 metastable level atoms,
	 eventually \Alert[red]{stimulated
	 emission rate will be
	 enhanced}
\end{itemize}
\end{minipage}
\vspace{2.5mm}


\xdef\loc{south east}
\begin{tikzpicture}[remember picture,overlay] % anchor=south
    \node (A) [yshift=-0.25*\figYshift,xshift=\figXshift, anchor=\loc] at (current page.\loc) {\includegraphics[page=1,width=9.75cm]{Figures_booklet/SOM_9_Statistical_mechanics_figures/SOM_Figure_26.png}}; %scale=\figscale. % 8cm max hei
\end{tikzpicture}


%\endofslide 
\endofsection
\end{frame}
%%%%%%%%%%%%%%%
